\section{Discussion}

\textbf{8.1 What transformation structure reveals about representations.} A representation carries more than what can be read from it: it carries what can still be enacted. Encoding merges states, and the mergings that survive transformation define an operable content with explicit conditions --- existence (compatibility), form (the measured organisation), use (the realised interface), and price (what composition costs). Asking what a network's features are becomes, in part, asking what they can still do to themselves. The reframing has a practical edge: an intervention that edits a feature vector is graded here by whether the \textit{remaining} layers produce what the real transformation produces, not by whether the edit looks large in feature space, and the two verdicts come apart routinely --- a random orthogonal operator tops the projection ranking at 15 of 42 sites while its transfer is negative at every one of them (\S{}5.5).

\textbf{8.2 What is general, and what is the colour instance.} The generality of the theory comes from its proofs, and it is worth stating precisely which parts are general. The fibre condition and the inherited composition law of \S{}3.2 hold for any encoder and any transformation family with a composition law; the risk decomposition of \S{}3.3 is an identity for any squared loss and any candidate family; the statement about coordinate-monotone families (Thm 3) is a restricted-class result about closed orbits; the accumulation law (Thm 4) needs only a Lipschitz truth action and a uniform single-step bound; the closure criterion (Props F--G) is stated for any carrier and action, which is why it can be applied to a measured organisation as a design condition; and carrier sufficiency (Thm 5) holds as an identity for any carrier, action and read-in, with the measured conditions sizing its terms, which is why it can return a \textit{negative} answer for an algebra --- as it does for the dissipative family. By contrast the kernel-preservation criterion and the closed form (Thms 1--2) are theorems about the single-region linear layer, and the harmonic concentration, the shared scaffold, the inheritance findings and the interface are colour \textit{measurements}: they are what this instance does, measured where the paper measures, and they are stated with the site and the grid attached. The boundary between the two is the claim--scope table of \S{}3.7, not the reader's inference.

\textbf{8.3 Where the framework stops.} Two limits bound the claims, and both are statements about the object rather than about the evidence for it. The first is the scope of the defect law: it is exact under isotropy, its tightness condition is the vanishing of a cross-covariance term (Proposition B.3), and the score depends on the input covariance (Corollary B.4). Real sites are strongly anisotropic, so the law is not a calibrated predictor there --- on 42 reachable sites the prediction is biased high by $+0.40$ to $+0.65$ and orders no better than the naive baselines --- and what it still supplies is the direction of the effect and a diagnostic use of its own deviation. The second is the dissipative family. The heat semigroup composes to $1.8\times10^{-7}$ at the truth level, yet its per-image decay rate is strongly content-dependent: the rate dispersion is $17.5$ and $8.6$ at layers 2 and 3 against $0.91$ for a synthetic carrier that is shared by construction, and a shared-rate predictor that beats the copy baseline $14.9\times$ on that carrier loses to it on the real sites ($0.83$ and $0.75$). \textbf{Figure 12} collects the out-of-sample tests of this section and of \S{}6.6: the closed form used as a predictor across the reachable sites, the sharing term inside and outside the fit split, the rank-resolved benefit of letting channels correct one another, and the shared-rate extrapolation itself. No basis we tested therefore supplies a common decay rate, the canonical spatial action does not compile, and the family is carried in channel space instead.

\begin{figure*}[t]
\centering
\includegraphics[width=\textwidth]{fig12_prediction.pdf}
\caption*{\textbf{Figure 12.} What a fitted law predicts out of sample. (a) Transfer predicted from each site's own mixing ratio by the closed form of Theorem 2 against the transfer measured at that site, one point per reachable site, with the identity line dotted and the measured mean dashed; the prediction is biased high by $0.40$ with mean absolute error $0.43$ against a measured mean of $0.31$, so across real sites the law orders no better than the naive baselines (Spearman $0.48$ between $\kappa$ and transfer, Pearson $-0.44$ between prediction and measurement) --- the negative result of \S{}8.3 shown against the line it would have to lie on. (b) Sharing term as a share of the realisation error, evaluated inside the fit split and on held-out data, per depth and algebra, on a logarithmic axis; the in-split estimate of $0.70$--$0.97$ collapses to $10^{-9}$--$0.073$ out of sample, so a dominant sharing cost is an artefact of fitting and evaluating in the same cells. (c) Error of the rank-$k$ subspace operator relative to the diagonal restriction, against depth, for three ranks and both algebras, with parity dotted; letting channels correct one another pays only at the first block (down to $0.25$ at rank $32$ under hue rotation) and costs at every deeper site, where the diagonal restriction predicts better. (d) Extrapolation to a diffusion time that was never fitted, from the shared mean rate of the fit times, as a gain over the copy baseline; the exact synthetic carrier reaches $14.9\times$ while the two real sites reach $0.83\times$ and $0.75\times$, and the number above each bar is that arm's rate dispersion ($0.91$ for the control against $17.5$ and $8.6$ at the two sites), which is why the shared-rate prediction does not transfer.}
\end{figure*}

\textbf{8.4 The criterion is the instrument, not the object.} The faithfulness score $T_F$ and its companions are properties of a \textit{pair} --- a representation and a consumer. They are not a scalar summary of a representation, and the paper's negative controls show how badly they can mislead when read that way: a random orthogonal operator attains the best projection coefficient at 15 of 42 sites, an unfaithful operator inflates the projection coefficient precisely by being unfaithful, and a consumer that does not read the transformed direction cannot grade any operator at all. The object of study is therefore the transformation's \textit{algebra} and the structure of its representation --- the induced action, its kernel and stabilisers, the four-block split, the gate-change term --- while the score is how we see it. This distinction also explains the paper's most easily misread result: the depth trend is \textit{not} a claim that deep layers represent fewer transformations. What declines with depth is the reach of fixed low-parameter families, and the decline tracks the readout--transformation pairing rather than the algebra (Table 11), which is the same statement as \S{}3.5's: a gate change counts only where the consumer reads it.

\textbf{8.5 Limits and threats to validity.} Six matter. (i) \textit{Residual domain gap.} The reference transformations are exact in rendering-parameter space, and the gap to real images is measured rather than assumed away: $3.4^\circ$ median read-out on synthetic shapes against $8.0^\circ$ on the high-concentration real regions, with the unaligned route at $69.5^\circ$. (ii) \textit{Consumer dependence.} Every capability number belongs to its consumer, which is why \S{}7.5 keeps the interface's matched-supervision comparison separate from the operator-family capacity controls. (iii) \textit{Single-seed cells.} The gate-change study, the semigroup-structure study and the mild-dissipation grid are single-seed, and single-seed cells are never pooled with three-seed cells. (iv) \textit{Estimator-limited fits.} Where the fitted dimension is large against the row count ($d^2/K=428$ at one site), the measured defect over-states the irreducible one, which is why every score carries its rank and budget. (v) \textit{Closure is a design criterion, not a diagnosis.} On measured features the synthesis has full column rank, the exact carrier model is closed automatically, and the residual failure of a global fit is attributable to cross-content inconsistency and tail rather than to non-closure; the criterion's role is to say which coordinate sets a compiled interface may keep. (vi) \textit{Partial algebra coverage.} Colour is carried through the full protocol; heat through composition, pairing and boundary arms; the remaining families are perimeter evidence, and the paper's claims about them are correspondingly narrower.

\textbf{8.6 The next explanatory step.} Three connections are open, and they are ordered by how much they would change the picture. The first is a \textit{nonlinear analogue of the mixing ratio}: where the site is differentiable, the exact intertwining condition implies the local condition $J\psi(\tau s)\,D\tau_s=D\rho_\tau(\psi(s))\,J\psi(s)$ on two Jacobians, so a locally defined $\kappa$, built from the local four-block decomposition, should predict the local defect --- which would turn the crossing decomposition of \S{}3.5 into a quantitative local law rather than a decomposition plus a correlation. The second is a \textit{carrier architecture}: layers that carry transformation laws explicitly, with defect budgets composed across depth the way \S{}5.4 composes them across a chain, would test whether the operability floor of \S{}6.1 is a property of learned representations or of how they are read. The third is the \textit{mechanism-level separation of learning's selection from its reorganisation} (\S{}4.5) --- currently the inheritance result shows training moves the organisation only modestly while the depth profile moves it a great deal, and the two effects are not yet separated by a controlled intervention. A resonance-based readout of shape--colour interaction remains a one-line speculation we do not develop here.
